\section{Conclusion}
The experiments separate adaptive value from detector choice and controller implementation. Primary Tabular-Q appears cheaper than Static-MedRF only because it fixes LightLR; under expanded load, the same frozen table falls to a TinyDT default with 67.76\% FPR, while DQN remains LightLR. In contrast, the simple causal DeadlineGuard switches between MedRF and LightLR, uses 43.7 J less than Static-MedRF (2.98\%), and removes all 512 measured deadline misses at a 1.373-pp F1 cost. The pairwise non-degeneracy test explains how reward normalization can suppress action reversal before deployment, and the off-support analysis shows how zero-initialized tabular ties can expose an unsafe default action. Learned IDS orchestration should therefore be credited with adaptive value only after realized actions, relevant static and conditional controls, service feasibility, security loss, and whole-device cost are separated.
